% !TEX root = ../main.tex

\section{NeoSetzer Feature Map}
\label{sec:feature-map}

This chapter is a map for readers who want to understand NeoSetzer beyond the
basic edit--build--preview loop. It is organized by task rather than by every
menu item. Settings, installed TeX tools, dictionaries, and keyboard shortcuts
can differ between systems, so use \menu{Preferences} and the Command Palette
as the live source of truth for your own installation.

\subsection{Find and run actions}
\label{sec:feature-map:actions}

Use the \textbf{Command Palette} when you know what you want to do but not
where the action lives. It groups commands by area, searches by name and
keywords, and explains when a command is unavailable in the current document
or view. This is often the fastest way to discover an operation before deciding
whether you want to assign or change a shortcut.

The welcome screen provides the other high-level entry points: create a blank
LaTeX or BibTeX document, open a file, start the Document Wizard, reopen recent
work, or create another writable example-project copy.

\subsection{Write efficiently}
\label{sec:feature-map:editing}

The editor combines syntax highlighting, automatic matching brackets,
auto-indentation, line operations, code folding, search and replace, and
context-aware completion. Completion adapts to LaTeX context: commands,
references, citations, labels, environments, symbols, enumerate items, and
Beamer frames are examples of suggestions that can be offered when the source
makes them meaningful. The \textbf{Shortcuts Bar} and \textbf{Symbols} sidebar
provide discoverable insertions when you prefer not to memorize commands.

Spell checking depends on an available dictionary and the language settings on
your system. It is best treated as a writing aid rather than a LaTeX validator:
use the build log for typesetting or package errors. Editor appearance, font,
line wrapping, indentation, current-line highlighting, and related behaviour
can be adjusted in the Appearance and Editor preference pages.

\subsection{Navigate, inspect, and organize}
\label{sec:feature-map:navigation}

The \textbf{Document Structure} sidebar lists headings, labels, project files,
and TODO entries. It shows ordinary section numbers when they are reliably
available from the source, including common appendices and literal counter
changes such as the sample in Appendix~\ref{sec:appendix-structure}. It does
not execute arbitrary TeX macros, so exotic number redefinitions may differ
from the final PDF.

Use the \textbf{Document Statistics} view when you need counts and reading
information, and use source labels and references to keep large documents
connected. The project and structure views are particularly useful with the
child files in this example project: opening a child does not prevent you from
building the root document.

\subsection{Build, diagnose, and review PDF output}
\label{sec:feature-map:build-preview}

NeoSetzer can use configured LaTeX build systems such as PDFLaTeX, XeLaTeX,
LuaLaTeX, Tectonic, and \texttt{latexmk} when they are installed. Save and
build is a common starting action, but shortcuts are configurable; inspect the
Keyboard Shortcuts preference page instead of relying on one fixed key.

The PDF preview supports zooming, searching within the PDF, and SyncTeX source
navigation when the selected engine produces the required synchronization data.
The build log groups diagnostics and lets you jump to reported source lines. If
a PDF changes outside NeoSetzer, the preview shows a persistent reload prompt
rather than silently leaving you with an old result.

\subsection{Keep projects safe and recoverable}
\label{sec:feature-map:recovery}

Auto-save can preserve recoverable editor state after a crash, while session
management and the recent-documents list help restore a working set. Recent
documents can be pinned. The Editor preferences also control how NeoSetzer
responds when source files change outside the application; enable automatic
reload only when that behaviour fits your workflow.

Auto-build can rebuild after editing stops, but it is deliberately a preference
rather than an assumption. It is most useful when the selected builder is fast
and the document is already in a healthy state.

\subsection{Tables, templates, and structured starting points}
\label{sec:feature-map:structured-tools}

The Document Wizard creates a new source document from configuration choices.
Its lightweight presets are separate from \textbf{user document templates},
which store complete LaTeX source snapshots for later preview, selection, and
reuse. The Insert Table dialog generates editable LaTeX grids, supports
constrained cell merges and longtable choices, and can import or paste UTF-8
CSV/TSV data. Chapter~\ref{sec:tables-data} contains a safe practice input.

\subsection{External tools and AI Fix}
\label{sec:feature-map:external-tools}

\textbf{AI Fix} is optional. When configured, it can hand a selected build
error and nearby source context to an external agent command. That agent runs
outside NeoSetzer and can modify files, so review its permissions, use trusted
project directories only, and keep ordinary backup or version-control habits.
Combining automatic external-file reload with auto-build can create a fast
workflow, but it also means external edits may be incorporated and rebuilt
quickly; enable that combination only when you understand the consequences.

\subsection{Where to tune behaviour}
\label{sec:feature-map:preferences}

\begin{description}[style=nextline]
    \item[Appearance]
        Select application and editor appearance, including themes and fonts.
    \item[Editor]
        Tune indentation, wrapping, matching behaviour, auto-save, external
        changes, and other writing-time behaviour.
    \item[Autocomplete]
        Choose completion behaviour and related suggestion settings.
    \item[Build System]
        Choose an installed builder and configure build-related behaviour such
        as automatic builds.
    \item[Keyboard Shortcuts]
        Inspect current bindings and change them to fit your workflow.
\end{description}

\subsection{A deliberate exploration path}
\label{sec:feature-map:explore}

For a complete first tour, use the Command Palette to find Document Structure,
Document Statistics, Symbols, Insert Table, Build Log, and Preferences. Then
return to this project, make one small change in each feature area, and build
again. The safest way to learn an editor is to connect each interface control
to a visible source or PDF result.
